\section{Demonstration: scoring a real monitor and a broken one}
\label{sec:demo}

A benchmark is only as convincing as what it does to a real monitor, so we score
two. The first is the companion composed timing protection level
\citep{baweja2026cti}---a genuine, published method---used as the \emph{reference
monitor}. The second is a deliberately under-bounded constant monitor, used as a
\emph{broken monitor} whose failures the benchmark should expose. Both are run
through the identical fault menu, at the identical alert limit and stated integrity
risk, by \code{examples/tib\_scorecard.rs}. Every number in this section is printed
by that example; none is transcribed by hand into the code, and all are
\Modelled{}---representative magnitudes, not operational tolerances, carrying no
certification implication.

\subsection{Setup}
\label{sec:demosetup}

The application alert limit is $\AL=\SI{100}{\nano\second}$---a demanding
time-transfer tolerance---and both monitors declare a stated integrity risk of
$\IR=\num{e-3}$. The reference monitor is the composed protection level
$\PL(\tau)=\max(\TPL_{\mathrm{handover}},\mathrm{HPL}(\tau))$ of the companion
paper, built from three representative sources (a GNSS-PPP source
$\{\sigma=\SI{2}{\nano\second},b=\SI{1}{\nano\second}\}$, a PTP/White-Rabbit source
$\{\SI{5}{\nano\second},\SI{2}{\nano\second}\}$, and a LEO-STL source
$\{\SI{8}{\nano\second},\SI{3}{\nano\second}\}$, each $p_{\mathrm{fault}}=\num{e-4}$)
and evaluated over a coast of one day discretized into $64$ epochs. It floors at
$\PL(0)=\SI{27.0}{\nano\second}$ at handover and rises to
$\SI{264.8}{\nano\second}$ at the $\SI{24}{\hour}$ end of the coast. The broken
monitor declares a constant $\PL=\SI{5}{\nano\second}$ at every epoch. The menu is
the nanosecond-scale stimulus of Table~\ref{tab:demomenu}, spanning the holdover
regime the reference monitor is designed for, the two undetectable delays, and
adversarial faults deliberately outside the holdover model.

\begin{table}[t]
  \centering
  \caption{The demonstration menu (\Modelled{}), $64$ epochs per scenario at
  $\AL=\SI{100}{\nano\second}$. Undetectable classes are fixed by Mizrahi
  \citep{mizrahi7384}.}
  \label{tab:demomenu}
  \begin{tabular}{l r l}
    \toprule
    scenario & peak offset & note \\
    \midrule
    domain divergence        & $\sim\!\SI{50}{\nano\second}$   & slow drift \\
    holdover coast           & (coast profile)                 & the regime the PL models \\
    static delay             & \SI{200}{\nano\second}          & constant injected delay \\
    incremental delay        & $\sim\!\SI{126}{\nano\second}$  & ramp \\
    clock slam               & \SI{150}{\nano\second}          & step at epoch 20 \\
    asymmetric delay         & \SI{40}{\nano\second}           & half of an \SI{80}{\nano\second} asymmetry \\
    path-selective           & \SI{50}{\nano\second}           & delay on half the paths \\
    $k$-of-$N$ quorum        & \SI{30}{\nano\second}           & 1 of 3 sources compromised \\
    symmetric delay          & \SI{20}{\nano\second}           & \textbf{undetectable} \\
    replay within freshness  & \SI{15}{\nano\second}           & \textbf{undetectable} \\
    \bottomrule
  \end{tabular}
\end{table}

\subsection{The reference monitor: covered where designed, honestly exposed elsewhere}
\label{sec:demoref}

Table~\ref{tab:refcard} is the reference monitor's scorecard. Three things happen,
and the benchmark reports all three without editorializing.

\begin{table}[t]
  \centering
  \caption{Scorecard of the reference monitor (the composed protection level,
  \SI{27.0}{\nano\second} at handover $\to$ \SI{264.8}{\nano\second} at
  \SI{24}{\hour}) at $\AL=\SI{100}{\nano\second}$, $\IR=\num{e-3}$. Total
  hazardous count $n_{\HMI{}}^{\text{total}}=64$;
  $\mathrm{all\_undetectable\_absorbed}=\text{true}$. \Modelled{}.}
  \label{tab:refcard}
  \begin{tabular}{l l r r r r}
    \toprule
    scenario & verdict & avail. & nom. & \MI{} & \HMI{} \\
    \midrule
    domain divergence        & \CovPass{}                    & \SI{66}{\percent} & \SI{66}{\percent} & 0  & 0  \\
    holdover coast           & \CovPass{}                    & \SI{66}{\percent} & \SI{66}{\percent} & 0  & 0  \\
    static delay             & \CovFail{}                    & \SI{66}{\percent} & \SI{0}{\percent}  & 0  & 42 \\
    incremental delay        & \CovPass{}                    & \SI{66}{\percent} & \SI{31}{\percent} & 22 & 0  \\
    clock slam               & \CovFail{}                    & \SI{66}{\percent} & \SI{31}{\percent} & 0  & 22 \\
    asymmetric delay         & \CovPass{}                    & \SI{66}{\percent} & \SI{23}{\percent} & 27 & 0  \\
    path-selective           & \CovPass{}                    & \SI{66}{\percent} & \SI{19}{\percent} & 30 & 0  \\
    $k$-of-$N$ quorum        & \CovPass{}                    & \SI{66}{\percent} & \SI{30}{\percent} & 23 & 0  \\
    symmetric delay          & \Absorbed{}   \emph{(undet.)} & \SI{66}{\percent} & \SI{66}{\percent} & 0  & 0  \\
    replay within freshness  & \Absorbed{}   \emph{(undet.)} & \SI{66}{\percent} & \SI{66}{\percent} & 0  & 0  \\
    \bottomrule
  \end{tabular}
\end{table}

\paragraph{It covers its design regime.} The two scenarios the composed protection
level is built for---slow domain divergence and free-running holdover coast---are
\CovPass{} with zero \HMI{}: the protection level bounds the growing coast error
throughout, and once the coast finally drives the protection level past the alert
limit the monitor declares \Unavail{}, a continuity event, never \HMI{}. This is
the intended behaviour: a coasting clock that outgrows its tolerance says so rather
than misleading.

\paragraph{Availability is uniform; the tightness spread lives in the nominal rate.}
Every scenario in Table~\ref{tab:refcard}---detectable and undetectable alike---has
the \emph{same} $\SI{66}{\percent}$ availability, and that uniformity is itself
informative. Availability here measures the monitor's self-declared usability
($1-n_{\Unavail{}}/n$), which is a property of the composed protection level
$\PL(\tau)$ against the alert limit, \emph{not} of the injected fault: the coast
crosses $\AL=\SI{100}{\nano\second}$ at the same epoch regardless of which fault is
present, so the same $\approx\!34\%$ of epochs are declared \Unavail{} on every
row. The entire per-fault differentiation therefore lives in the two columns that
\emph{do} depend on the fault: the \emph{nominal rate}, which falls from
$\SI{66}{\percent}$ on the clean coast to $\SI{19}{\percent}$ under path-selective
delay as the protection level, though still available, grows loose relative to a
larger true error (a tightness signal), and the \HMI{} count, which stays zero
except where an adversarial fault drives the true error past the alert limit. Reading
availability and the nominal rate as separate columns is what keeps continuity and
tightness from being conflated.

\paragraph{It absorbs both undetectables.} The symmetric delay
($\SI{20}{\nano\second}$) and the replay ($\SI{15}{\nano\second}$) both sit under
the $\SI{27}{\nano\second}$ handover floor, so $\PL(k)\ge|e(k)|$ at every epoch and
both earn \Absorbed{}; $\mathrm{all\_undetectable\_absorbed}$ is true. The monitor
is credited for having \emph{reserved enough protection level} to carry faults it
could never detect---the only honest defence (Section~\ref{sec:absorptionrule})---and
is nowhere credited for the impossible act of detecting them.

\paragraph{It is honestly exposed at its boundary.} The benchmark does not flatter
the reference monitor. Adversarial faults \emph{outside} the holdover model produce
hazardously misleading information: the static $\SI{200}{\nano\second}$ delay draws
$42$ \HMI{} epochs and \CovFail{}, and the $\SI{150}{\nano\second}$ clock slam
draws $22$ \HMI{} and \CovFail{}---in both cases the true error exceeds the
$\SI{100}{\nano\second}$ alert limit while the protection level, tuned for holdover,
does not cover it. The mid-range faults (incremental, asymmetric, path-selective,
quorum) pass coverage but accrue \MI{}: the protection level under-bounds them
\emph{within} the alert limit, a tightness signal, not a hazard. The reference
monitor's total hazardous count is $n_{\HMI{}}^{\text{total}}=64$. This is the
demonstration's thesis in one number: a reference yardstick \emph{maps} a method's
coverage---it confirms the composed protection level is sound in its holdover
regime and reveals precisely where it needs the adversarial extensions of the later
papers in the arc---rather than declaring it uniformly safe.

\subsection{The broken monitor: strictly worse, on the same rule}
\label{sec:demobroken}

Table~\ref{tab:brokencard} scores the deliberately under-bounded
$\SI{5}{\nano\second}$ monitor. It is worse everywhere the rule can see. Its total
hazardous count is $n_{\HMI{}}^{\text{total}}=121$, nearly double the reference's;
the static delay alone draws $64$ \HMI{} (every epoch hazardous, since even the
$\SI{5}{\nano\second}$ bound plus the $\SI{100}{\nano\second}$ limit cannot contain
a $\SI{200}{\nano\second}$ error), the clock slam draws $44$, and the incremental
delay now pierces the alert limit for $13$ \HMI{}. Decisively, \emph{both}
undetectable offsets are now \Unabsorbed{}: a $\SI{5}{\nano\second}$ protection
level cannot carry a $\SI{20}{\nano\second}$ or $\SI{15}{\nano\second}$ offset, so
$\mathrm{all\_undetectable\_absorbed}$ is false. Yet even here---this is the
invariant of Section~\ref{sec:verdict} at work---neither undetectable scenario is
ever reported as detected; the honest verdict for a monitor that under-reserves
against an undetectable fault is \Unabsorbed{}, not ``detection missed.''

\paragraph{Availability $\neq$ integrity, made concrete.} The sharpest reading of
Table~\ref{tab:brokencard} is its availability column: the broken monitor is
\SI{100}{\percent} \emph{available} on \emph{every} scenario. Its
$\SI{5}{\nano\second}$ protection level never exceeds the $\SI{100}{\nano\second}$
alert limit, so it never once declares itself \Unavail{}---by the Stanford--WAAS
definition it is maximally available. Yet the same monitor is riddled with
misleading and hazardously misleading epochs: $121$ \HMI{} in total, and \MI{} on
nearly every otherwise-``available'' epoch. This is the concrete statement that
availability is not integrity. A monitor that never says ``unavailable'' but returns
hazardously wrong answers is not safer than one that honestly withdraws---it is
worse, because it invites reliance precisely where reliance is unwarranted. The
reference monitor's lower $\SI{66}{\percent}$ availability, read naively, looks like
the weaker number; read correctly, it is the honest one, because the epochs it
declined were epochs on which its bound could no longer be trusted.

\begin{table}[t]
  \centering
  \caption{Scorecard of the broken monitor (constant \SI{5}{\nano\second}) on the
  same menu, alert limit, and stated risk. Total hazardous count
  $n_{\HMI{}}^{\text{total}}=121$;
  $\mathrm{all\_undetectable\_absorbed}=\text{false}$. \Modelled{}.}
  \label{tab:brokencard}
  \begin{tabular}{l l r r r r}
    \toprule
    scenario & verdict & avail. & nom. & \MI{} & \HMI{} \\
    \midrule
    domain divergence        & \CovPass{}                    & \SI{100}{\percent} & \SI{11}{\percent} & 57 & 0  \\
    holdover coast           & \CovPass{}                    & \SI{100}{\percent} & \SI{19}{\percent} & 52 & 0  \\
    static delay             & \CovFail{}                    & \SI{100}{\percent} & \SI{0}{\percent}  & 0  & 64 \\
    incremental delay        & \CovFail{}                    & \SI{100}{\percent} & \SI{5}{\percent}  & 48 & 13 \\
    clock slam               & \CovFail{}                    & \SI{100}{\percent} & \SI{31}{\percent} & 0  & 44 \\
    asymmetric delay         & \CovPass{}                    & \SI{100}{\percent} & \SI{0}{\percent}  & 64 & 0  \\
    path-selective           & \CovPass{}                    & \SI{100}{\percent} & \SI{0}{\percent}  & 64 & 0  \\
    $k$-of-$N$ quorum        & \CovPass{}                    & \SI{100}{\percent} & \SI{0}{\percent}  & 64 & 0  \\
    symmetric delay          & \Unabsorbed{} \emph{(undet.)} & \SI{100}{\percent} & \SI{0}{\percent}  & 64 & 0  \\
    replay within freshness  & \Unabsorbed{} \emph{(undet.)} & \SI{100}{\percent} & \SI{0}{\percent}  & 64 & 0  \\
    \bottomrule
  \end{tabular}
\end{table}

\subsection{What the demonstration shows, and what it does not}
\label{sec:demoscope}

The demonstration establishes that the benchmark discriminates a sound monitor from
a broken one on a common rule, rewards honest reservation against the undetectable,
and refuses to credit detection of the undetectable in either case. What it does
\emph{not} establish is any accuracy claim: it does not certify that the composed
protection level is safe in absolute terms, that $\SI{27}{\nano\second}$ is the
right floor, or that these fault magnitudes are operationally representative---the
magnitudes are \Modelled{} illustration. Nor can a $64$-epoch run resolve a
deep-tail \num{e-7} integrity rate; the stated risk here is $\num{e-3}$, the rate a
short menu can support. The demonstration's role is to show the yardstick working,
honestly, on a real method---mapping its coverage rather than flattering it---and
that is exactly what an honesty-immune benchmark should do.
